% Integration fragment; requires amsmath, amssymb, amsthm and natbib.
% Environments theorem, proposition, corollary, remark are supplied by the main file.
% Verified primary references: chen2022, zhang,
% bounemoura, delshams1996, poschel.

\subsection{A statistical bridge in the norm required by KAM theory}
\label{sec:statistical-bridge}

All Hamiltonians below are real-valued scalar functions, including when their
parameters and hidden representations are complex. The realification of a complex
parameter vector has dimension $P$. Let $\mathcal D=\mathbb T^n\times B$,
where $B$ is a bounded ball in action space, and write $d=2n$.
Lebesgue measure on $\mathcal D$ is denoted by $dx$.
The following result completes the derivative-norm interface needed to combine
the learning argument of \citet{chen2022} with a finite-differentiability KAM
theorem. It does not infer derivative smallness from function values alone.

\begin{theorem}[Uniform learning-to-$C^\ell$ bound]
\label{thm:learning-to-cl}
Fix $d<p<\infty$, an integer $\ell\geq1$, and an integer $s>\ell+d/p$.
Let $X_1,\ldots,X_N$ be independent samples with density
$f_\mu\geq m_0>0$ on $\mathcal D$.
Let $H_\star$ be the target Hamiltonian and
$\{H_\vartheta:\vartheta\in\Theta\}$ the hypothesis class, with
$\Theta\subset[-B_\Theta,B_\Theta]^P$.
Assume that for all $x\in\mathcal D$ and $\vartheta,\eta\in\Theta$,
\begin{align}
 \|\nabla H_\vartheta(x)-\nabla H_\star(x)\|_p&\leq Q,\label{eq:residual-bound}\\
 \|\nabla H_\vartheta(x)-\nabla H_\eta(x)\|_p
   &\leq L_g\|\vartheta-\eta\|_\infty.\label{eq:parameter-bound}
\end{align}
After fixing the additive gauge by
\begin{align}
 e_\vartheta&=H_\vartheta-H_\star-
   \frac{1}{|\mathcal D|}\int_{\mathcal D}(H_\vartheta-H_\star)\,dx,
 \nonumber\\
 \sup_{\vartheta\in\Theta}\|e_\vartheta\|_{W^{s,p}(\mathcal D)}&\leq M_s,
 \label{eq:uniform-sobolev}
\end{align}
suppose that $M_s$ is a valid uniform bound. Define
\begin{align}
 \widehat R_N(\vartheta)
 &=\frac1N\sum_{i=1}^N
    \|\nabla H_\vartheta(X_i)-\nabla H_\star(X_i)\|_p^p,\\
 L_{\mathrm{par}}&=pQ^{p-1}L_g,\qquad L_{\max}=Q^p.
\end{align}
For any fixed $0<a\leq2B_\Theta$ and $0<\delta<1$, put
\begin{align}
 \Gamma_N(a,\delta)
 &=2L_{\mathrm{par}}a\nonumber\\
 &\quad+\frac{L_{\max}}{\sqrt{2N}}
 \sqrt{P\log(1+2B_\Theta/a)+\log(1/\delta)},\\
 S_N(\vartheta)&=\widehat R_N(\vartheta)+\Gamma_N(a,\delta),\\
 G_N(\vartheta)&=C_{\mathrm P}\,m_0^{-1/p}S_N(\vartheta)^{1/p},\\
 \alpha&=\frac{\ell-1+d/p}{s-1}\in(0,1).
\end{align}
Here $C_{\mathrm P}$ is a Poincare constant chosen so that a mean-zero function
satisfies $\|e\|_{W^{1,p}}\leq C_{\mathrm P}\|\nabla e\|_{L^p}$.
With probability at least $1-\delta$, simultaneously for every
$\vartheta\in\Theta$,
\begin{equation}
 \|e_\vartheta\|_{C^\ell(\overline{\mathcal D})}
 \leq E_N(\vartheta):=
 C_{\mathrm{GN}}\left[
 G_N(\vartheta)+M_s^\alpha G_N(\vartheta)^{1-\alpha}\right].
 \label{eq:learning-cl-bound}
\end{equation}
$C_{\mathrm{GN}}$ depends only on the domain and $(d,p,s,\ell)$.
We use $\|e\|_{C^\ell}=\max_{|a|\leq\ell}\sup|D^ae|$;
equivalent conventions only change the displayed constants.
The event is uniform over the hypothesis class and therefore includes parameters
selected using these same training data.
\end{theorem}

\begin{proof}
The loss $r_\vartheta(x)=\|\nabla H_\vartheta(x)-\nabla H_\star(x)\|_p^p$
lies in $[0,L_{\max}]$ and is $L_{\mathrm{par}}$-Lipschitz in the
$\ell^\infty$ parameter metric. Partitioning the parameter cube into cells of
side length $a$ and choosing a representative from each nonempty cell gives
an $a$-net with at most $(1+2B_\Theta/a)^P$ points.
The one-sided Hoeffding inequality and a union bound control
$R(\eta)-\widehat R_N(\eta)$ at every net point by the square-root term in
$\Gamma_N$, where $R(\eta)=\int r_\eta\,d\mu$.
The population and empirical losses each change by at most
$L_{\mathrm{par}}a$ between a parameter and its net representative.
Consequently $R(\vartheta)\leq S_N(\vartheta)$ uniformly.

The density lower bound gives
$\|\nabla e_\vartheta\|_{L^p(dx)}\leq m_0^{-1/p}S_N(\vartheta)^{1/p}$.
The gauge in \eqref{eq:uniform-sobolev} and Poincare's inequality then yield
$\|e_\vartheta\|_{W^{1,p}}\leq G_N(\vartheta)$.
Apply the Gagliardo--Nirenberg inequality to $g=\nabla e_\vartheta$,
with derivative orders $j=\ell-1$, $m=s-1$, and both endpoint
Lebesgue exponents equal to $p$. Its interpolation identity is
\[
 0=\frac{j}{d}+\alpha\left(\frac1p-\frac{m}{d}\right)
     +\frac{1-\alpha}{p},
\]
which gives the stated $\alpha$.
The strict inequality $s>\ell+d/p$ ensures $\alpha<1$ and avoids the
critical Sobolev endpoint. Extension from the smooth bounded action domain,
together with periodic charts in the angular variables, gives
\[
 \|e\|_{C^\ell}
 \leq C_{\mathrm{GN}}\left(
 \|e\|_{W^{1,p}}+
 \|e\|_{W^{1,p}}^{1-\alpha}\|e\|_{W^{s,p}}^\alpha\right).
\]
Lower-order derivatives are included by the inhomogeneous term and the
$W^{1,p}\hookrightarrow C^0$ embedding, which uses $p>d$.
Substitution proves \eqref{eq:learning-cl-bound}.
\end{proof}

For instance, one may select an integer $\ell>2(\tau+1)$, $p>2n$, and
$s=\ell+1$. This makes every exponent explicit, although the resulting
interpolation bound can be conservative. Bounds on $M_s,Q,L_g$ follow from
bounded parameters and derivatives of smooth activations on a compact domain;
the corresponding bounds for $H_\star$ must also be justified. Merely declaring
the functions smooth does not provide these quantitative bounds.
Alternatively, a loss containing all spatial derivatives of the vector field
through order $\ell$ controls $W^{\ell+1,p}$ directly, and Sobolev embedding
then gives a $C^\ell$ error with exponent $1/p$, without interpolation.

\begin{corollary}[KAM transfer for the target and the trained Hamiltonian]
\label{cor:common-target-kam}
Let $n\geq2$, $\tau>n-1$ and an integer $\ell>2(\tau+1)$ be fixed.
Suppose that, in one specified canonical action--angle chart,
\[
 H_\star=h_0(I)+P_\star(I,\varphi),\qquad
 H_\vartheta=h_0(I)+P_\vartheta(I,\varphi).
\]
Assume $h_0\in C^{\ell+2}(\overline B)$ with finite norm, and
$\nabla h_0:B\to\Omega$ is a diffeomorphism with bounded inverse in $C^\ell$.
Fix $\gamma>0$ and set
\begin{equation}
 \Omega_{\gamma,\tau}=\left\{\omega\in\Omega:\begin{array}{l}
 |k\cdot\omega|\geq\gamma |k|_1^{-\tau}\ (0\ne k\in\mathbb Z^n),\\
 \operatorname{dist}(\omega,\partial\Omega)\geq\gamma
 \end{array}\right\}.
\end{equation}
Let $c_{\mathrm K}>0$ be the constant in Theorem A of
\citet{bounemoura} for this domain and integrable Hamiltonian.
If a certified bound $\beta_\vartheta\geq
\|P_\vartheta-\overline P_\vartheta\|_{C^\ell}$ satisfies
\begin{equation}
 \beta_\vartheta+E_N(\vartheta)<c_{\mathrm K}\gamma^2,
 \label{eq:joint-kam-budget}
\end{equation}
then, on the event in Theorem~\ref{thm:learning-to-cl}, both $H_\star$ and
$H_\vartheta$ possess invariant Lagrangian tori for every frequency in
$\Omega_{\gamma,\tau}$. The flow on each torus is $C^1$-conjugate to the
linear flow with that frequency. For a piecewise smooth frequency-domain
boundary, the complement of each family has measure at most $C\gamma$.
In particular, each family has positive measure when $\gamma$ is small enough.
\end{corollary}

\begin{proof}
Write $\overline P$ for the spatial mean over $\mathcal D$.
The gauge-fixed difference satisfies
\[
 P_\star-\overline P_\star
   =P_\vartheta-\overline P_\vartheta-e_\vartheta.
\]
Thus the two perturbation norms are at most $\beta_\vartheta$ and
$\beta_\vartheta+E_N(\vartheta)$, respectively. Constants added to a
Hamiltonian do not change its vector field. Apply the cited finite-smoothness
KAM theorem separately to the two Hamiltonians with the same $h_0$,
domain, and Diophantine set. The theorem supplies tori of class
$C^{\tau+1}$ with Lipschitz dependence as invariant graphs on frequency.
The two systems need not have identical invariant embeddings.
\end{proof}

The certified budget includes the physical perturbation from $h_0$, any
normal-form remainder, and the learning error. A small learning error alone
does not establish that the physical system is close to an integrable one.
Theorem A in the published 2020 version of \citet{bounemoura} gives
the measure bound $C\gamma$; the $C\gamma^2$ measure bound printed in the
2018 preprint should not be used.

\subsection{A coefficient certificate for a finite analytic CHNN}
\label{sec:coefficient-certificate}

The finite analytic readout used in the experiments yields a stronger certificate
that avoids a dimension-dependent interpolation loss. Let $\psi_1,\ldots,
\psi_m$ be real-valued basis functions on the real phase domain. Complex
features may be converted into this basis by taking their real and imaginary
parts and collecting equal monomials. The certificate applies to the effective
coefficients of the actual trained Hamiltonian, irrespective of how they are
parameterized inside the network.

\begin{proposition}[Deterministic global certificate from sampled gradients]
\label{prop:gram-certificate}
Suppose
\begin{align*}
 H_\vartheta&=H_{\mathrm{prior}}+\sum_{j=1}^m c_{\vartheta,j}\psi_j,\\
 H_\star&=H_{\mathrm{prior}}+\sum_{j=1}^m c_{\star,j}\psi_j+R.
\end{align*}
The basis contains no constant gauge direction. Let
\begin{align*}
 G_N&=\frac1N\sum_{i=1}^N D\psi(X_i)^\top D\psi(X_i),\\
 \lambda_N&=\lambda_{\min}(G_N)>0,
\end{align*}
where the $j$th column of $D\psi$ is $\nabla\psi_j$.
For gradient observations $Y_i=\nabla H_\star(X_i)+\xi_i$, define
\begin{align*}
 r_N&=\left(\frac1N\sum_i\|\nabla H_\vartheta(X_i)-Y_i\|_2^2\right)^{1/2},\\
 t_N&=\left(\frac1N\sum_i\|\nabla R(X_i)+\xi_i\|_2^2\right)^{1/2}.
\end{align*}
Then
\begin{equation}
 \|c_\vartheta-c_\star\|_2\leq\frac{r_N+t_N}{\sqrt{\lambda_N}}.
 \label{eq:coefficient-error}
\end{equation}
For any function norm $\|\cdot\|_{\mathcal X}$ satisfying the triangle
inequality, put $K_{\mathcal X}=(\sum_j\|\psi_j\|_{\mathcal X}^2)^{1/2}$.
If $\|R\|_{\mathcal X}\leq\varepsilon_R$, then
\begin{equation}
 \|H_\vartheta-H_\star\|_{\mathcal X}
 \leq K_{\mathcal X}\frac{r_N+t_N}{\sqrt{\lambda_N}}+\varepsilon_R.
 \label{eq:global-coefficient-certificate}
\end{equation}
The norm may be $C^\ell$ on a real compact domain or the supremum norm on
a prescribed complex neighborhood where every basis function and $R$ is
holomorphic. No independent sampling assumption is needed.
\end{proposition}

\begin{proof}
Let $\Delta c=c_\vartheta-c_\star$.
The empirical norm of $D\psi\Delta c$ is
$(\Delta c^\top G_N\Delta c)^{1/2}$.
Since $D\psi\Delta c=(\nabla H_\vartheta-Y)+(\nabla R+\xi)$ at every
sample, the triangle inequality gives
$\sqrt{\lambda_N}\|\Delta c\|_2\leq r_N+t_N$.
Finally,
$\|\sum_j\Delta c_j\psi_j-R\|_{\mathcal X}
\leq K_{\mathcal X}\|\Delta c\|_2+\varepsilon_R$ by Cauchy--Schwarz.
\end{proof}

For raw monomials $\psi_\nu(x)=x^\nu$, an explicit $C^\ell$ bound on
$[-r,r]^d$ is obtained from
\[
 \|\psi_\nu\|_{C^\ell}
 \leq\max_{|a|\leq\ell,\,a\leq\nu}
      \frac{\nu!}{(\nu-a)!}\,r^{|\nu|-|a|}.
\]
The analytic sup norm on $\{|x_j|\leq r\}\subset\mathbb C^d$ is
$r^{|\nu|}$. For all 35 degree-four monomials in four variables,
$K_{\mathrm{an}}\leq\sqrt{35}\,r^4$.
If a canonical chart $x=\Phi(I,\varphi)$ maps an action--angle complex
domain into this polydisc, the same analytic bound holds after composition.
For a $C^\ell$ certificate after composition, the derivative bounds of
$\Phi$ must instead be included by the chain rule.
Unmodeled physical terms require the tail bounds $t_N$ and $\varepsilon_R$.
Their approximation-tail contribution is zero only when the target lies
exactly in the declared class. The full-cosine feature class below has
this property for its specified reduced target; it does not remove
truncation error relative to the full beam.

The statistical result requires independent state samples with the stated
coverage density. Consecutive states on a trajectory do not satisfy this
assumption automatically. Accordingly, the trajectory-based experiments
below use the deterministic Gram certificate, which requires no sampling
independence. Their implementation divides both squared residuals and the
Gram matrix by $4N$ rather than $N$; the coefficient bound is unchanged
provided the noise norm uses that same scalar-observation normalization.

